\chapter{Realization}

\section{introduction} 
This chapter is dedicated to the practical realization of the IO‑Box project, focusing on
the firmware development carried out across the different software layers. While the
previous chapters introduced the system concept, architecture, and components, the
present chapter illustrates the concrete implementation work that transformed the design
into a functional prototype.

The realization process is organized around three main layers:
\begin{itemize}
    \item \textbf{Board Support Package (BSP)}: Development of low‑level drivers and
    routines specific to the STM32G474 microcontroller and its peripherals, ensuring
    proper initialization and hardware control.
    \item \textbf{Manager Layer}: Implementation of supervisory modules responsible
    for protocol handling, error detection, and data integrity, providing a reliable
    interface between BSP and application logic.
    \item \textbf{Application Layer}: Integration of functional logic, including signal
    acquisition, filtering, regulation, actuator control, and communication with external
    supervisory systems.
\end{itemize}

By detailing the coding and integration performed in each of these layers, this chapter
demonstrates how the IO‑Box firmware was realized in practice, highlighting the modular
approach and ensuring traceability from design specifications to executable code.

\section{Firmware Architecture}


The firmware of the IO‑Box was designed according to a layered architecture, ensuring
modularity, portability, and clear separation of responsibilities. Figure~\ref{fig:fw_architecture}
illustrates the overall structure, which is organized into six main components:

- \textbf{Application layer}: Peripheral calls and business logic. Handles high‑level
commands, events, and callbacks.

- \textbf{Manager layer}: Supervises peripherals and protocols (EEPROM, UART, RS485,
SPI/I2C sensors, DAC, IO logic, CAN‑FD). Provides state machines, buffering,
coordination, and bus/channel mapping.

- \textbf{BSP layer}: Implements low‑level drivers (UART, SPI, I2C, ADC, PWM, CAN,
USB, DAC, etc.) as HAL wrappers, exposing standardized APIs to higher layers.

- \textbf{Utilities}: Shared helpers such as \texttt{Utils.h}, board configuration, validation
macros, and debug/logging.

- \textbf{Callbacks \& ISR}: HAL interrupts, callback functions, and event posting.

- \textbf{HAL layer}: Vendor‑provided hardware abstraction library used internally by BSP.

\begin{figure}[H]
    \centering
    \includegraphics[width=0.85\textwidth]{C:/Users/selim/Downloads/imagesreport/fw_architecture.png}
    \caption{Layered firmware architecture of the IO‑Box}
    \label{fig:fw_architecture}
\end{figure}

This architecture provides a clear hierarchy: the BSP abstracts hardware access, the
Manager coordinates protocols and resources, and the Application layer implements the
functional logic. Utilities and callbacks support all layers, while HAL remains the base
hardware abstraction.


\section{Board Support Package (BSP)}

The Board Support Package (BSP) represents the lowest software layer of the IO‑Box
firmware. It provides the essential drivers and initialization routines that allow the
STM32G474 microcontroller to interact with its peripherals in a standardized way.
By abstracting hardware details, the BSP ensures that higher layers of the firmware
can access resources such as timers, ADCs, communication interfaces, and memory
modules without depending directly on register‑level configurations.

In the context of this project, the BSP was developed to replace generic HAL usage
with custom, instance‑specific drivers. This approach guarantees modularity,
portability, and precise control over hardware behavior, forming the foundation
upon which the Manager and Application layers are built.

\subsection{CubeMX Project Configuration}
This subsection presents the global configuration of the STM32G474 microcontroller
performed in CubeMX, including pin assignments and clock setup. Figure~\ref{fig:pinout}
illustrates the complete pinout used in the IO‑Box firmware project.

\begin{figure}[H]
    \centering
    \includegraphics[width=0.75\textwidth]{C:/Users/selim/Downloads/imagesreport/stm32_pinout.png}
    \caption{STM32G474 CubeMX pinout configuration}
    \label{fig:pinout}
\end{figure}

\subsubsection{SPI2 BSP Implementation}

The SPI2 BSP was implemented to enable communication with the EEPROM memory
module integrated by the hardware team. Configuration parameters such as baud rate,
polarity, phase, and data size were defined in \texttt{BSP\_SPI2.h} using macros. Figure~\ref{fig:spi2_config}
shows the CubeMX configuration interface used to initialize SPI2 in Mode~0 with a
prescaler of 16.

\begin{figure}[H]
    \centering
    \includegraphics[width=0.75\textwidth]{C:/Users/selim/Downloads/imagesreport/spi_macro.png}
    \caption{example of SPI2 configuration in SPI2BSP.h}
    \label{fig:spi2_config}
\end{figure}

The API follows the same structure applied to all BSP drivers:
- \textbf{Initialization}: \texttt{BSP\_SPI2\_init}, \texttt{BSP\_SPI2\_pins\_init}, \texttt{BSP\_SPI2\_dma\_init}
- \textbf{Basic I/O}: \texttt{BSP\_SPI2\_read}, \texttt{BSP\_SPI2\_write}, \texttt{BSP\_SPI2\_transfer}
- \textbf{Non‑blocking I/O}: \texttt{BSP\_SPI2\_read\_IT}, \texttt{BSP\_SPI2\_write\_DMA}, etc.
- \textbf{Control}: Chip Select management, reset/deinit, buffer flush
- \textbf{State}: \texttt{BSP\_SPI2\_getstate} for monitoring peripheral status

This structured design ensures that SPI2 can reliably serve the EEPROM module while
remaining reusable for other peripherals if needed. The detailed EEPROM management
logic will be presented in the Manager section.
 
\begin{figure}[H]
    \centering
    \includegraphics[width=0.75\textwidth]{C:/Users/selim/Downloads/imagesreport/spi_API.png}
    \caption{SPI2 API's in SPI2BSP.h}
    \label{fig:spi2_API}
\end{figure}